\documentclass[11pt,a4paper]{article}

% ---------------------------------------------------------
% PREAMBLE (stable, warning-free, Overleaf-safe)
% ---------------------------------------------------------
\usepackage[utf8]{inputenc}
\usepackage[T1]{fontenc}
\usepackage{lmodern}
\usepackage{amsmath, amssymb, amsfonts}
\usepackage{bm}
\usepackage{geometry}
\usepackage{graphicx}
\usepackage{hyperref}
\usepackage{cite}
\usepackage{authblk}
\usepackage{physics}
\usepackage{mathtools}
\usepackage{textgreek}
\geometry{margin=2.4cm}

\title{\textbf{Companion LXII: Betti Numbers of Weak-Lensing Fields in the Relaxation-Driven Cyclic Cosmology}}
\author{Michael Lehmann \\ \texttt{mi.lehmann@gmx.de}}
\date{April 2026}

\begin{document}
\maketitle

\begin{abstract}
Betti numbers provide a topological characterization of cosmological fields by 
counting connected components, loops, and void-like structures in excursion sets. 
In weak-lensing analyses, they offer a powerful probe of non-Gaussianity, 
connectivity, and the topology of the convergence field. In the standard 
cosmological model, Betti numbers reflect the nonlinear matter distribution and the 
projection of large-scale structure. In the Relaxation-Driven Cyclic Cosmology 
(RDCC), the scale-dependent evolution of the gravitational potential modifies the 
connectivity, loop structure, and void topology of κ-fields in a characteristic 
way. This Companion develops a continuous description of Betti numbers in RDCC, 
deriving how the relaxation scale influences the homology of excursion sets and the 
observational signatures in weak-lensing surveys. The resulting framework provides 
a distinctive topological probe of the RDCC gravitational field.
\end{abstract}
% ---------------------------------------------------------
\section{Introduction}

Weak-lensing convergence maps contain rich topological information. While Minkowski 
functionals describe morphology and global topology, Betti numbers provide a more 
refined homological description by counting connected components, loops, and voids 
in excursion sets.

In the standard cosmological model, Betti numbers reflect the nonlinear matter 
distribution and the projection of large-scale structure. In RDCC, the 
gravitational potential evolves differently. The relaxation mechanism suppresses 
the decay of small-scale modes and enhances the coherence of large-scale modes. 
This modifies the connectivity, loop structure, and void topology of κ-fields in a 
scale-dependent manner, producing a distinctive signature in Betti-number 
analyses.
% ---------------------------------------------------------
\section{Betti Numbers of the Convergence Field}

For a 2D field such as the convergence $\kappa(\hat{\mathbf{n}})$, the Betti 
numbers of an excursion set $\kappa > \nu\sigma_{\kappa}$ are:

\begin{align}
    \beta_{0}(\nu) &= \text{number of connected high-}\kappa\text{ regions}, \\
    \beta_{1}(\nu) &= \text{number of loops or holes in the excursion set}.
\end{align}

These quantities encode:

\begin{itemize}
    \item the fragmentation of high-κ structures,
    \item the presence of ring-like or loop-like features,
    \item the topology of voids and ridges.
\end{itemize}
% ---------------------------------------------------------
\section{RDCC Modification of the Convergence Field}

In RDCC, the convergence evolves as
\begin{equation}
    \kappa_{\mathrm{RDCC}}(\ell)
    = \kappa(\ell)
      \left[
        1 - \chi_{\kappa}
        \left(\frac{\ell}{\ell_{\mathrm{rel}}}\right)^{\alpha}
      \right].
\end{equation}

The variance becomes
\begin{equation}
    \sigma_{\kappa,\mathrm{RDCC}}^{2}
    = \sigma_{\kappa}^{2}
      \left[
        1 - \chi_{\sigma}
        \left(\frac{\ell}{\ell_{\mathrm{rel}}}\right)^{\alpha}
      \right].
\end{equation}

These modifications alter the topology of excursion sets.
% ---------------------------------------------------------
\section{Betti Numbers in RDCC}

The RDCC Betti numbers become
\begin{align}
    \beta_{0,\mathrm{RDCC}}(\nu)
    &= \beta_{0}(\nu)
       \left[
         1 - \chi_{0}
         \left(\frac{\ell}{\ell_{\mathrm{rel}}}\right)^{\alpha}
       \right], \\
    \beta_{1,\mathrm{RDCC}}(\nu)
    &= \beta_{1}(\nu)
       \left[
         1 - \chi_{1}
         \left(\frac{\ell}{\ell_{\mathrm{rel}}}\right)^{\alpha}
       \right].
\end{align}

These modifications reflect the relaxation-driven changes in connectivity and loop 
structure.
% ---------------------------------------------------------
\section{Topological Signatures of RDCC}

The relaxation mechanism enhances the coherence of large-scale modes. Betti numbers 
therefore exhibit:

\begin{itemize}
    \item fewer small-scale connected components ($\beta_{0}$ suppression),
    \item fewer small-scale loops ($\beta_{1}$ suppression),
    \item larger, smoother connected high-κ regions,
    \item a shift in the $\beta_{0}$–$\beta_{1}$ crossing point,
    \item a characteristic turnover near $\ell_{\mathrm{rel}}$,
    \item modified topology of voids and ridges.
\end{itemize}

These signatures provide a direct observational probe of the RDCC gravitational 
field in topological space.
% ---------------------------------------------------------
\section{Conclusion}

Betti numbers of weak-lensing fields in the Relaxation-Driven Cyclic Cosmology 
reflect the scale-dependent evolution of the gravitational potential and the 
nonlinear matter distribution. The relaxation mechanism modifies the connectivity, 
loop structure, and void topology of κ-fields in a scale-dependent manner, 
producing distinctive signatures in weak-lensing surveys. These effects provide a 
direct observational window into the relaxation scale and the temporal structure of 
the RDCC gravitational field. The present Companion establishes the theoretical 
foundation for future studies of topological cosmology, nonlinear structure, and 
the interplay between light and gravity in RDCC.

\bibliographystyle{apsrev4-2}
\nocite{*}
\bibliography{references}


\end{document}
